\documentclass[12pt,a4paper]{article}
\usepackage[utf8]{inputenc}
\usepackage{geometry}
\usepackage{fancyhdr}
\usepackage{hyperref}

% Geometría de la página
\geometry{
 a4paper,
 total={170mm,257mm},
 left=20mm,
 top=20mm,
}

% Encabezado y pie de página
\pagestyle{fancy}
\fancyhf{}
\fancyhead[C]{Propuesta de Webinar: Seguridad en Criptoactivos}
\fancyfoot[C]{\thepage}
\renewcommand{\headrulewidth}{0.4pt}
\renewcommand{\footrulewidth}{0.4pt}

% Título
\title{
    \vspace{2cm}
    \textbf{Propuesta de Webinar Gratuito} \\
    \large \textit{Blockchain: La Revolución de la Confianza Digital y sus Beneficios Frente al Sistema Financiero Tradicional}
    \vspace{1.5cm}
}
\author{Ing. Daniel Velarde Kubber}
\date{La Paz, 13 de junio de 2025}

\begin{document}

\maketitle
\thispagestyle{empty}
\clearpage

\section*{1. Detalle General de la Propuesta}

\begin{itemize}
    \item \textbf{Título del Webinar:} Blockchain: La Revolución de la Confianza Digital y sus Beneficios Frente al Sistema Financiero Tradicional.
    \item \textbf{Propuesta dirigida a:} Universidad Real de la Cámara Nacional de Comercio.
    \item \textbf{Ponente:} Ing. Daniel Velarde Kubber - Ingeniero Industrial y de Sistemas, entusiasta en Data y Blockchain.
    \item \textbf{Fecha propuesta:} 13 de junio de 2025.
    \item \textbf{Ciudad:} La Paz, Bolivia.
    \item \textbf{Modalidad:} Virtual (a través de la plataforma que disponga la universidad).
    \item \textbf{Duración Total:} 1 hora (45 minutos de exposición y 15 minutos para preguntas y respuestas).
    \item \textbf{Costo:} Gratuito.
\end{itemize}

\section*{2. Objetivo del Webinar}

El webinar busca explicar de manera clara y concisa la tecnología Blockchain, la base sobre la que se construyen los criptoactivos. El objetivo es que los participantes comprendan sus principios fundamentales, como la descentralización y la inmutabilidad, y analicen los beneficios disruptivos que ofrece en comparación con la arquitectura del sistema financiero y administrativo tradicional, enfatizando siempre la seguridad y la implementación responsable.

\section*{3. Temario y Contenido (Duración: 45 minutos)}

\begin{enumerate}
    \item \textbf{El Problema de la Confianza y los Intermediarios (5 min)}
    \begin{itemize}
        \item Cómo funciona el sistema tradicional: Bancos, notarios, gobiernos como garantes de confianza.
        \item Puntos débiles: Centralización, costos, lentitud y puntos únicos de fallo.
    \end{itemize}
    
    \item \textbf{¿Qué es Blockchain? Una Analogía Comprensible (15 min)}
    \begin{itemize}
        \item Blockchain como un "libro de contabilidad digital, público e inmutable".
        \item Componentes clave:
        \begin{itemize}
            \item \textbf{Bloques:} Contenedores de transacciones.
            \item \textbf{Cadena:} El enlace criptográfico que garantiza la integridad.
            \item \textbf{Red Descentralizada (P2P):} Cómo se distribuye y valida la información sin una autoridad central.
        \end{itemize}
        \item Tipos de Blockchain: Públicas (Bitcoin, Ethereum) vs. Privadas (Hyperledger).
    \end{itemize}
    
    \item \textbf{Beneficios Comparados con el Sistema Financiero Existente (15 min)}
    \begin{itemize}
        \item \textbf{Transparencia vs. Opacidad:} Registros públicos y auditables en tiempo real.
        \item \textbf{Seguridad vs. Vulnerabilidad:} Eliminación de puntos únicos de fallo y resistencia a la manipulación de datos.
        \item \textbf{Eficiencia vs. Burocracia:} Reducción de intermediarios, costos y tiempos de liquidación (ej. Contratos Inteligentes).
        \item \textbf{Inclusión Financiera vs. Barreras de Acceso:} Potencial para bancarizar a la población no atendida.
    \end{itemize}
    
    \item \textbf{Seguridad, Responsabilidad y Casos de Uso (10 min)}
    \begin{itemize}
        \item La seguridad intrínseca de la tecnología vs. la seguridad en su aplicación.
        \item Riesgos: Ataques del 51\%, vulnerabilidades en contratos inteligentes, y la importancia de la gestión de claves.
        \item Más allá de las finanzas: Aplicaciones en cadenas de suministro, gestión de identidad, votaciones y más.
        \item La necesidad de un desarrollo y adopción responsables.
    \end{itemize}
\end{enumerate}

\section*{4. Sesión de Preguntas y Respuestas (15 minutos)}

Al finalizar la presentación, se abrirá un espacio para que los asistentes puedan realizar preguntas y resolver sus dudas directamente con el ponente.

\vspace{2cm}

\begin{center}
    Agradezco de antemano su tiempo y consideración. \\
    \vspace{0.5cm}
    \textbf{Ing. Daniel Velarde Kubber} \\
    \textit{Entusiasta en Data y Blockchain}
\end{center}

\end{document}